%% Introduction - Le pourquoi du comment on dit ça %%

Le projet de cartographie du traitement documentaire réalisé à l'INA nécessite avant tout de comprendre les raisons de son existence afin de déterminer à quels besoins il doit répondre. En effet, ces collections sont au contact de plusieurs typologies de publics avec leurs propres pratiques, usages et besoins. Ils accèdent aux collections des institutions patrimoniales par des outils différents et ont en conséquent une expérience différente d'une typologie de public à l'autre. C'est pourquoi nous allons dans un premier temps poser les bases de notre réflexion en analysant ces enjeux de l'accessibilité aux collections patrimoniales. 

%% Plan %%

Dans ce sens, nous allons décrire et analyser les publics des institutions patrimoniales en les divisant en catégories de publics : le grand public, les chercheurs, les documentalistes et les professionnels de l’audiovisuel. Nous décrirons leurs habitudes, leurs pratiques et les outils à leurs dispositions avant de préciser, dans une seconde partie les différents freins à l'accessibilité qu'ils rencontrent. 